\icode{N-MCH-CA-L2-260627-001}
将扁平的石子向水面快速抛出，石子可能会在水面上一跳一跳地飞向远方，俗称“打水漂”。要使石子从水面跳起产生“水漂”效果，石子接触水面时的速度方向与水面的夹角不能大于\(\theta\)。为了观察到“水漂”，一同学将一石子从距水面高度为\(h\)处水平抛出，抛出速度的最小值为多少?（不计石子在空中飞行时的空气阻力，重力加速度大小为\(g\)）